\section{Retained singular dynamics and interface conventions}\label{sec:appendix}
\subsection{The nonreset singular realization}
The technical companion~\cite{companion}, Section~5, supplies a second direct verification of the block hypotheses, distinct from Gaussian filtering. We record its raw data and the exact point of attachment so that its status is not confused with the adaptive progressive-sensor class of Section~\ref{sec:adaptive}.

Take ratios $r_l\in[1/5,1/3]$, $\ell_l=\prod_{j\le l}r_j$ and a fair binary Cantor coordinate $x=\sum_l(1-r_l)\ell_{l-1}\zeta_l$. Complete its domain by adjoining every finite-prefix fair-tail mean. A prepared mode $I\in\{0,1\}$ is reported, followed by $B$ paid initial prefix reads. At a later opportunity an independent active indicator has any deterministic probability. Inactivity is an exact copy but costs a call. On activity, mode zero either prepends two fair reported bits, deletes the first bit within mode zero, or changes to mode one while deleting; the probabilities are $(1-\alpha_t)(1-\gamma_0)$, $(1-\alpha_t)\gamma_0$, and $\alpha_t$, with $\gamma_0=1/4096$. Mode one prepends six reported fair bits and moves to zero with probability $1/2$, or copies within mode one with probability $1/2$. Deletion does not reveal the deleted bit; both mode changes transport a nonconstant old tail.

Within each mode give the completed coordinate Euclidean distance weighted by $v_0=1800$, $v_1=1$; between modes use distance $2\sqrt{1800}$. The relation is the full same-interface candidate domain and $w=1$. On the completed domain deletion has gain at most $30$ and prepending $s$ fixed bits has gain at most $6\,3^{-s}$. The matrix of physical expected squared gains is
\[
 C_t=\begin{pmatrix}(1-\alpha_t)85/128&900\alpha_t\\2/59049&1/2\end{pmatrix},
 \qquad C_t\binom{1800}{1}\le\frac{17}{20}\binom{1800}{1}.
\]
Cross-mode pairs have squared gain at most $1/4$. Thus the active block certificate is $m=1$, $\kappa=17/20$, and constant envelope. All finite-prefix centers through depth $h$ in both modes form a global Borel cover of cardinality $2^{h+2}-2$ and radius $\sqrt{1800}\ell_h$. The gain estimates follow by comparing the first distinguishing level: its original separation lies between $\ell_{k-1}/6$ and $\ell_{k-1}$, while deletion or prepending shifts this level. These facts verify the hypotheses independently of any claimed memory bound.

The companion proves the resulting exact finite-acquisition law, including the physical-oracle lower and the mode-dependent affine terminal Bernoulli task. If $H_T$ is the actual acquired depth, then the actual score law is a \emph{finite atomic} mixture of the fair-prefix means. Conditional tail variance is comparable to $\E\ell_{H_T}^2$, and a physical Cantor small-ball witness gives the memory term $\ell_{\lfloor\log_2M\rfloor}^2$. Their maximum and sum are both comparable to
\[
 \E\ell_{\min(H_T,\lfloor\log_2M\rfloor)}^2.
\]
The matched sign-minimax law adds $\delta^2$ and charges $B+T+2$ raw calls. The physical law is used solely as an oracle lower witness, never substituted for the actual atomic acquired law. Its proof and all additional assumptions are retained in full in the companion. The progressive acquisition theorem here is a different result: it permits adaptive allocation across coordinates and even smooth factors, but does not replace this nonreset dynamical theorem.

For completeness, exact independent holds preserve the active error radius without requiring a uniform positive activation probability. If active updates satisfy $e_+\le\sqrt\kappa e+\Delta$ and holds copy exactly with independent probability $1-\vartheta_t$, then
\[
 e_t^2\le(1-\vartheta_t)e_{t-1}^2+
       \vartheta_t(\sqrt\kappa e_{t-1}+\Delta)^2.
\]
Both terms preserve $e\le\max\{e_0,\Delta/(1-\sqrt\kappa)\}$. The envelope argument is identical with indicator-form expectations. No conditioning on a null active event is needed when $\vartheta_t=0$. State-dependent missingness in Section~\ref{sec:filters} is not such an independent hold.

\subsection{Finite interfaces and the location of cuts}
A publicly chosen terminal index after an encoding cut means that each retained label selects a vector, or a response function, of possible readouts. It does not mean that the encoder saw the selected index before that cut. This is why the continuation function in Example~\ref{ex:query} has dimension $d$ while its terminal answer has only two values. When a clock or an interface tag can be queried after the cut, its information is part of the retained tuple unless the theorem explicitly conditions on that public interface.

An exact continuation challenge can force disjoint state supports. For a tag $J$ with probabilities $p_j$, unreissued after the cut, exact recovery forces at least one disjoint label set of size $m_j$ per tag. If the score's conditional law is $\nu_j$, the precise distortion is
\[
 \inf_{m_j\ge1,\ \sum_jm_j\le K}\sum_jp_jq_{m_j}(\nu_j).
\]
The proof fixes independent coins, observes that one label cannot decode two distinct tags exactly, and applies conditional nearest-center quantization within each support. Disjoint codebooks give the reverse inequality. This inherited tag mechanism is different from the noisy continuation-function bound of Lemma~\ref{lem:cut}; neither licenses an unchallenged simulator-state lower bound.

On the analytic side all report versions used by wrong candidates must be jointly Borel on the declared true/candidate domain. In the Gaussian case emission densities are positive and the update is everywhere defined on the open simplex. In the continuous-instrument proposition only the positive-density domain is normalized. The root allowance $u$ in Theorem~\ref{thm:transfer} is valid only for routines still satisfying the relation, not for arbitrary floating-point perturbations of the same size. A calibrated-input perturbation can be included in $u$ only after proving that same stable-metric and relation certificate.

These conventions also delimit the retained companion statements. Finite-dimensional identities there do not imply unbounded-operator closure; finite stopped total variation does not imply a large-deviation principle; and same-controller coupling does not imply unrestricted optimal control. The independent realization articles on ordered measurements, finite physical actions and streaming space are not premises of the foundation theorem.
